\documentclass[11pt]{article}
\usepackage[margin=1in]{geometry}
\usepackage{booktabs}
\usepackage{amsmath}
\usepackage[hidelinks]{hyperref}
\title{Clustering Related Messages in LLM Chat Sessions for Guided Agent Search}
\author{Mash project --- experimental working draft}
\date{October 5, 2026}
\begin{document}
\maketitle
\begin{abstract}
We investigate flat, overlapping groups of related messages within agent chat
threads, using existing embeddings without generative model inference. A first
bounded comparison of four clustering families on 24 real threads found that
context-aware average linkage recovered distant returning subjects but did not
satisfy preregistered quality gates. We document this negative result and a
separate expanded screen of five additional families, none of which passed its
development screen. This is a living research
record, not a completed publication: references are sparse and agent-authored,
and planned experiments are explicitly distinguished from measurements.
\end{abstract}

\section{Task and motivation}
An agent searching a long thread needs messages addressing the same concrete
subject, even when intervening turns address other subjects. Shared repository,
vocabulary, or broad intent alone does not establish relatedness. A mixed-subject
message may belong to several groups; noise may remain unassigned. The output
contains source identities and actual representative messages, not generated
summaries. Global topics and a user-facing hierarchy are outside scope.
Contiguous dialogue segmentation motivates local coherence checks
\cite{hearst,xing}, but cannot itself recover distant returning subjects.

\section{Representation and scoring}
The existing local index supplies 384-dimensional all-MiniLM-L6-v2 message
embeddings, produced by quantized CPU ONNX inference and stored as symmetric
signed bytes. Decoding validates dimensions and stored norms. Experiments use
normalized float64 vectors, without new inference, PCA, training, or an LLM.
Canonical history is unchanged. Exact prepared projections are deduplicated for
fitting; all original source identities remain in the scoring map.

The second representation in the first study fixes context weights:
\[
 z_i=\operatorname{normalize}(0.7v_i+0.3v_{p(i)}).
\]
An explicit earlier parent is preferred; an absent parent uses the immediately
preceding visible embedded message if eligible. Explicit roots and missing
explicit parents abstain. Eligibility is a structural proxy: at least five
whitespace tokens, excluding whole injected instruction/envelope headers.
This proxy is not semantic verification and never reads reference labels.

Groups need two eligible unique projections. Final mean centroids supply
memberships above the configured cosine cutoff and within 0.05 of the best
centroid score; this step does not change the fitted partition. Representatives
are observed messages closest to those means. We score primary partitions and
overlap separately. Only explicitly judged pairs enter quality metrics;
ambiguous pairs are excluded, and unjudged pairs are never negative labels.
Precision, recall, and F1 are macro-averaged per thread. Coverage counts reviewed
substantive messages, not every message passing the structural proxy.

\section{First study: completed}
\subsection{Corpus and experimental procedure}
The read-only Xenia inventory contained 42,150 message vectors in 817 threads,
from Claude Code and Codex. Twenty-four independent components were sampled:
12 development and 12 held out, two threads per provider and short/medium/long
stratum per split. Length bins were 8--24, 25--150, and 151+ vectors. All 817
source files were checked for known lineage. Provider ancestry, common sessions,
and conservative exact-projection copy overlap kept related components together.
Development contained 11,081 vectors; held out contained 4,705.

Before candidate outputs, 282 agent-authored pair judgments were frozen:
development 71 related, 57 unrelated, 14 ambiguous; held out 70, 59, 11.
Long messages were reviewed as excerpts, not independently annotated in full.
Twelve settings compared online cosine leader, BIRCH, cosine average linkage,
and HDBSCAN. Another twelve used the fixed context view. One setting per family
was chosen by development F1; held-out data were scored once. No settings were
tuned on held-out failures. Exact settings, hashes and codec contracts are
recorded in \path{docs/message-clustering-evaluation.md} and the locked Python
experiment. BIRCH used individual partial fits and reassigned all messages to
the final leaf centres; HDBSCAN noise was never one shared cluster.

\subsection{Acceptance criteria}
Promotion required macro precision $\geq .85$, recall $\geq .70$, F1 $\geq .75$,
reviewed substantive coverage $\geq .80$, and F1 gain $\geq .10$ over the best
simple baseline. Returning-subject, short-reply and mixed-message precision
needed to reach .75 with adequate controls. No provider/representation stratum
could lose more than .10 F1 against its best baseline. At least 90\% of audited
groups needed coherent representatives and farthest members. Warm grouping
p95 limits were 2 seconds at 1,000 messages and 15 seconds at 5,000, with process
peak RSS at most 512 MiB. Determinism and source identity preservation under
append/edit/delete/restart were additional contracts. Failure after the fixed
budget meant rejecting that pipeline, not indefinitely extending its grid.

\subsection{Quality results}
\begin{center}
\begin{tabular}{lrrrr}
\toprule
Held-out method & Precision & Recall & F1 & Coverage\\
\midrule
Leader + context & .719 & .842 & .736 & .945\\
BIRCH + context & .842 & .749 & .763 & .963\\
Average linkage + context & .839 & .815 & .774 & .963\\
HDBSCAN + context & .687 & .573 & .618 & .862\\
Whole thread & .543 & 1.000 & .698 & 1.000\\
Chronological windows (8) & .839 & .741 & .769 & .991\\
\bottomrule
\end{tabular}
\end{center}
Every method failed precision and the required baseline gain. Average linkage
recovered 11/13 returning-subject positives versus 3/13 for windows, but its
overall F1 gain was only .005. All held-out case-tagged pairs were positive;
apparent case precision lacked negative controls and could not pass that gate.
A fixed-hash sample audited 11 average-link groups from four held-out threads,
using representatives and two farthest primary members: nine were coherent.
Failures grouped unrelated Git progress or distinct session-related problems.
This suggests a representation/context issue worth testing, but does not prove
it is the dominant bottleneck or that other algorithms cannot help.

\subsection{Cost and update results}
Each method ran in a separate Python process on Xenia with one numerical thread.
Ten warm rebuilds estimated sample p95; three controlled prepared-unit updates
per operation measured rebuild cost. Model inference, preparation, transport,
cache lookup and cold imports were excluded. On one real 3,775-unit thread,
average-link warm p95 was 1.678 seconds and peak RSS 327 MiB; append/edit/delete
rebuild medians were 1.351/1.325/1.328 seconds. At 5,000 synthetic units its
p95 was 2.476 seconds and RSS 431 MiB. Synthetic scaling repeated real vectors
with seeded tiny jitter and provides no quality evidence. HDBSCAN exceeded the
memory gate at 5,000 (700 MiB); all methods met grouping-time gates. Repeated
result hashes matched within each controlled operation. These are small-sample
Python measurements, not production Rust or end-to-end latency claims.

\section{Expanded algorithm screen: completed development comparison}
The user requested both additional offline candidates and a larger cohort.
We froze 72 fresh components: 48 development, 24 reserved. We excluded all components
containing a first-study thread. The available Claude medium and long strata
have only six fresh independent candidates each; use Claude short/medium/long
allocation 18/6/6 and Codex 18/12/12, with two-thirds development in every stratum.
This availability-limited cohort is not balanced or representative of provider
market share. Selection seed and SHA order are fixed before candidate outputs;
long strata prioritize threads above 1,000 vectors.

Keep representation, retention, and cosine overlap cutoff .35 fixed. Compare
complete linkage (cuts .35/.50/.65), Euclidean DP-means \cite{dpmeans}
(squared penalties 1.30/1.00/.70), OPTICS (xi .03/.05/.10), spectral clustering
(4/8/16 groups), and affinity propagation (preference quantiles .25/.50/.75).
Implementation parameters are fixed in
\path{docs/message-clustering-alternatives.md}. Each fit uses an isolated
single-threaded Linux process with a 30-second deadline. Nonconvergence, failures,
timeouts, and peak RSS are reported; incomplete candidates cannot be ranked as
complete successes. Sparse development references must be frozen before fit
outputs. Reserved threads remain unjudged and unscored until a candidate passes
development screening. No new settings are added after outputs.

The expanded corpus contains 12,037 vectors: 5,498 development and 6,539 reserved.
The longest development thread has 716 messages; the reserved split includes a
4,118-message thread. Thus this screen does not establish long-thread scaling.
Each development thread contributes six sampled positions (two at the start,
midpoint, and end) and six fixed adjacent/cross-block pair judgments. Review used
140-character heads and 90-character tails. The 288 agent-authored references
contain 111 related, 117 unrelated, and 60 ambiguous pairs. These sparse excerpts
are diagnostic, not independent human ground truth. Exact duplicate projections
were checked; no case-specific quality claim follows from these labels.

\begin{center}
\begin{tabular}{lrrr}
\toprule
Best complete development setting & Precision & Recall & F1\\
\midrule
Complete-link (.35 cut) & .773 & .701 & .696\\
DP-means (.70 squared penalty) & .684 & .833 & .716\\
OPTICS (.03 xi) & .650 & .875 & .710\\
Spectral (8 clusters) & .755 & .748 & .716\\
Affinity propagation (.25 quantile) & .703 & .661 & .652\\
Average-link reference & .720 & .819 & .716\\
Chronological windows (8) & .773 & .748 & .718\\
Whole thread & .517 & .979 & .653\\
\bottomrule
\end{tabular}
\end{center}
None reached development precision .85, F1 .75, or a .10 baseline gain. Across
768 fits including the reference, 766 completed; affinity quantiles .50 and .75
each failed convergence on one thread. Partial-cohort metrics were not ranked as
complete results. No timeouts or other errors occurred. The fixed screen therefore
stopped without judging or scoring the 24 reserved threads. These development
scores cannot be compared directly to the first study's held-out scores.

Maximum observed completed-fit times across settings were .078 seconds
(complete-link), 1.133 (DP-means), 1.226 (OPTICS), .045 (spectral), .247
(affinity), and .063 (average-link). Corresponding process peak RSS maxima were
164/163/175/190/197/166 MiB, including inherited forked input state. Each fit was
measured once; these are not p95 estimates or end-to-end latency. Complete settings,
failed runs, source hashes and per-provider metrics are recorded in the companion
protocol/results document. Average-link primary-only F1 (.728) exceeded its overlap
F1 (.716), whereas OPTICS/spectral benefited substantially from overlap. This
motivates isolating the shared membership heuristic in subsequent work, rather
than attributing all quality differences to the partition algorithm.

GRINCH \cite{grinch} is relevant as an incremental hierarchy with rotations and
grafting, followed by a flat cut and separate overlapping assignment. It is not
implemented or evaluated in this batch. Online insertion is not required:
cached on-request affected-thread recomputation is eligible, but no such cache
or production clustering feature has been implemented.

\section{Next representation experiment: proposed, not measured}
The user proposed clustering prompt/answer pairs rather than individual messages.
Question and answer may be complementary, particularly for short referential
turns, but long answers may cover multiple subjects. Compare prompt-only units,
jointly embedded prompt/answer text, and prompt-conditioned bounded answer
passages against the existing message pipeline. Pooled existing vectors provide
a separate cheap control; they are not equivalent to jointly embedding text.
Token limits and truncation, multiple assistant messages, tool/progress filtering,
explicit parent boundaries and active incomplete turns must be specified before
outputs. Preserve original source identities and evaluate the same judged message
pairs, not a denominator changed by the new unit definition. This proposal has
no measurements yet and does not reopen the completed fifteen-setting screen.

\section{Limitations and reproducibility}
Only two providers, sparse agent-authored excerpt judgments, a structural
support proxy, and one timing host limit interpretation. Human independent
annotation, case-specific negative controls, broader provider coverage and
fresh held-out evaluation remain necessary. Similar-thread search and thread
vector quality are follow-ons. Source histories, embeddings, detailed judgments
and results remain outside Git in private directories; paper tables contain
aggregate measurements only. The first study's negative decision does not
establish that message clustering is impossible.

\begin{thebibliography}{9}
\bibitem{hearst} M. Hearst. TextTiling: Segmenting Text into Multi-paragraph
Subtopic Passages. Computational Linguistics, 1997.
\url{https://aclanthology.org/J97-1003/}
\bibitem{xing} L. Xing and G. Carenini. Improving Unsupervised Dialogue Topic
Segmentation with Utterance-Pair Coherence Scoring. SIGDIAL, 2021.
\url{https://aclanthology.org/2021.sigdial-1.18/}
\bibitem{dpmeans} B. Kulis and M. Jordan. Revisiting k-means: New Algorithms
via Bayesian Nonparametrics. ICML, 2012.
\url{https://icml.cc/2012/papers/291.pdf}
\bibitem{grinch} N. Monath et al. Scalable Hierarchical Clustering with Tree
Grafting. KDD, 2019. \url{https://arxiv.org/abs/2001.00076}
\end{thebibliography}
\end{document}
